% GENERATED FILE. Edit canonical lessons, then run npm run generate:books.
% canonical-source-hash: fnv1a64:9a67733657fdcb12
% canonical-lessons: ML-C265-potiyuka, ML-C265-kattikkuka, ML-C265-ketuttuka, ML-C265-kettuka, ML-C265-alakkuka

\chapter{To wrap, To light, To put out, To tie, To measure}
\label{ch:ml-to-wrap-to-light-to-put-out-to-tie-to-measure}
\hlchaptermodality{265}

\begin{chapteropening}
\textbf{By the end of this chapter:} \emph{I can say to wrap, to light, to put out, to tie and to measure.}

Complete the last lesson of chapter 265: I can say to wrap, to light, to put out, to tie and to measure.
\end{chapteropening}

\section[\texorpdfstring{\ml{പൊതിയുക} (potiyuka)}{പൊതിയുക (potiyuka)}]{\ml{പൊതിയുക} (potiyuka) — to wrap}

\label{lesson:ML-C265-potiyuka}



\begin{quote}
Before the new one: say the Malayalam for to sweep, then the Malayalam for to fold.
\end{quote}

\subsection*{You'll want to know: \ml{പൊതിയുക}}
\textbf{\ml{പൊതിയുക}} — \emph{potiyuka} — \textquotedblleft{}to wrap\textquotedblright{}.

\begin{cousinweb}
\textbf{\ml{പൊതിയുക}} (\emph{potiyuka}) is native Dravidian, one word with Tamil \emph{putai}, Kannada \emph{pode} and Telugu \emph{podugu} (DEDR 4509).
\end{cousinweb}

\begin{grammarlens}[title={What you've built}]
One more word for this chapter.
\end{grammarlens}

\subsection*{Your turn}
\begin{itemize}
\raggedright
  \item \emph{Say it:} \emph{potiyuka}
  \item \emph{Say it:} \emph{potiyuka}, once more
  \item \emph{Say it:} say \emph{potiyuka}
  \item \emph{Recall:} say the Malayalam for to sweep, then the Malayalam for to fold, then say \emph{potiyuka} again
\end{itemize}

\begin{tcolorbox}[breakable,colback=teal!4,colframe=teal!35!black,title={Before you move on}]
What does \ml{പൊതിയുക} mean? (\textquotedblleft{}To wrap\textquotedblright{}.) Say it once more.
\end{tcolorbox}
\section[\texorpdfstring{\ml{കത്തിക്കുക} (kattikkuka)}{കത്തിക്കുക (kattikkuka)}]{\ml{കത്തിക്കുക} (kattikkuka) — to light, to burn}

\label{lesson:ML-C265-kattikkuka}



\begin{quote}
Before the new one: say the Malayalam for to fold, then the Malayalam for to wrap.
\end{quote}

\subsection*{You'll want to know: \ml{കത്തിക്കുക}}
\textbf{\ml{കത്തിക്കുക}} — \emph{kattikkuka} — \textquotedblleft{}to light, to burn\textquotedblright{}.

\begin{cousinweb}
\textbf{\ml{കത്തിക്കുക}} (\emph{kattikkuka}) is native Dravidian, the causative of the verb \textquotedblleft{}to kindle, burn\textquotedblright{}; DEDR 1207 lists the root with Kannada \emph{kattu}.
\end{cousinweb}

\begin{grammarlens}[title={What you've built}]
One more word for this chapter.
\end{grammarlens}

\subsection*{Your turn}
\begin{itemize}
\raggedright
  \item \emph{Say it:} \emph{kattikkuka}
  \item \emph{Say it:} \emph{kattikkuka}, once more
  \item \emph{Say it:} say \emph{kattikkuka}
  \item \emph{Recall:} say the Malayalam for to fold, then the Malayalam for to wrap, then say \emph{kattikkuka} again
\end{itemize}

\begin{tcolorbox}[breakable,colback=teal!4,colframe=teal!35!black,title={Before you move on}]
What does \ml{കത്തിക്കുക} mean? (\textquotedblleft{}To light, to burn\textquotedblright{}.) Say it once more.
\end{tcolorbox}
\section[\texorpdfstring{\ml{കെടുത്തുക} (keṭuttuka)}{കെടുത്തുക (keṭuttuka)}]{\ml{കെടുത്തുക} (keṭuttuka) — to put out, to switch off}

\label{lesson:ML-C265-ketuttuka}



\begin{quote}
Before the new one: say the Malayalam for to wrap, then the Malayalam for to light.
\end{quote}

\subsection*{You'll want to know: \ml{കെടുത്തുക}}
\textbf{\ml{കെടുത്തുക}} — \emph{keṭuttuka} — \textquotedblleft{}to put out, to switch off\textquotedblright{}.

\begin{cousinweb}
\textbf{\ml{കെടുത്തുക}} (\emph{keṭuttuka}) is native Dravidian, the causative of the verb \textquotedblleft{}to be extinguished, be spoilt\textquotedblright{}; DEDR 1942 lists the root with Tamil \emph{keṭu} and Kannada \emph{keḍu}.
\end{cousinweb}

\begin{grammarlens}[title={What you've built}]
One more word for this chapter.
\end{grammarlens}

\subsection*{Your turn}
\begin{itemize}
\raggedright
  \item \emph{Say it:} \emph{keṭuttuka}
  \item \emph{Say it:} \emph{keṭuttuka}, once more
  \item \emph{Say it:} say \emph{keṭuttuka}
  \item \emph{Recall:} say the Malayalam for to wrap, then the Malayalam for to light, then say \emph{keṭuttuka} again
\end{itemize}

\begin{tcolorbox}[breakable,colback=teal!4,colframe=teal!35!black,title={Before you move on}]
What does \ml{കെടുത്തുക} mean? (\textquotedblleft{}To put out, to switch off\textquotedblright{}.) Say it once more.
\end{tcolorbox}
\section[\texorpdfstring{\ml{കെട്ടുക} (keṭṭuka)}{കെട്ടുക (keṭṭuka)}]{\ml{കെട്ടുക} (keṭṭuka) — to tie}

\label{lesson:ML-C265-kettuka}



\begin{quote}
Before the new one: say the Malayalam for to light, then the Malayalam for to put out.
\end{quote}

\subsection*{You'll want to know: \ml{കെട്ടുക}}
\textbf{\ml{കെട്ടുക}} — \emph{keṭṭuka} — \textquotedblleft{}to tie\textquotedblright{}.

\begin{cousinweb}
\textbf{\ml{കെട്ടുക}} (\emph{keṭṭuka}) is native Dravidian, one word with Tamil \emph{kaṭṭu}, Kannada \emph{kaṭṭu} and Telugu \emph{kaṭṭu} (DEDR 1147). DEDR 1147 compares Sanskrit \emph{ghaṭṭa}.
\end{cousinweb}

\begin{grammarlens}[title={What you've built}]
One more word for this chapter.
\end{grammarlens}

\subsection*{Your turn}
\begin{itemize}
\raggedright
  \item \emph{Say it:} \emph{keṭṭuka}
  \item \emph{Say it:} \emph{keṭṭuka}, once more
  \item \emph{Say it:} say \emph{keṭṭuka}
  \item \emph{Recall:} say the Malayalam for to light, then the Malayalam for to put out, then say \emph{keṭṭuka} again
\end{itemize}

\begin{tcolorbox}[breakable,colback=teal!4,colframe=teal!35!black,title={Before you move on}]
What does \ml{കെട്ടുക} mean? (\textquotedblleft{}To tie\textquotedblright{}.) Say it once more.
\end{tcolorbox}
\section[\texorpdfstring{\ml{അളക്കുക} (aḷakkuka)}{അളക്കുക (aḷakkuka)}]{\ml{അളക്കുക} (aḷakkuka) — to measure}

\label{lesson:ML-C265-alakkuka}



\begin{quote}
Before the new one: say the Malayalam for to put out, then the Malayalam for to tie.
\end{quote}

\subsection*{You'll want to know: \ml{അളക്കുക}}
\textbf{\ml{അളക്കുക}} — \emph{aḷakkuka} — \textquotedblleft{}to measure\textquotedblright{}.

\begin{cousinweb}
\textbf{\ml{അളക്കുക}} (\emph{aḷakkuka}) is native Dravidian, one word with Tamil \emph{aḷa}, Kannada \emph{aḷe} and Telugu \emph{alavi} (DEDR 295).
\end{cousinweb}

\begin{grammarlens}[title={What you've built}]
One more word for this chapter.
\end{grammarlens}

\subsection*{Your turn}
\begin{itemize}
\raggedright
  \item \emph{Say it:} \emph{aḷakkuka}
  \item \emph{Say it:} \emph{aḷakkuka}, once more
  \item \emph{Say it:} say \emph{aḷakkuka}
  \item \emph{Recall:} say the Malayalam for to put out, then the Malayalam for to tie, then say \emph{aḷakkuka} again
\end{itemize}

\begin{tcolorbox}[breakable,colback=teal!4,colframe=teal!35!black,title={Before you move on}]
What does \ml{അളക്കുക} mean? (\textquotedblleft{}To measure\textquotedblright{}.) Say it once more.
\end{tcolorbox}
